\chapter{Message Queuing and Streaming Ingestion with Amazon Kinesis}
\index{Amazon SQS}\index{Amazon SNS}\index{Kinesis Data Streams}\index{shard}\index{partition key}\index{Amazon MSK}\index{Apache Kafka}%
\begin{objectives}
\begin{itemize}
    \item Decouple producers from consumers with Amazon SQS queues and SNS fan-out patterns.
    \item Size Kinesis Data Streams from shard throughput limits and design partition keys for per-entity ordering.
    \item Write retrying Boto3 producers and idempotent consumers, and interpret \texttt{IteratorAgeMilliseconds}.
    \item Compare Kinesis with Amazon MSK (managed Kafka) on operations, pricing, and ecosystem fit.
    \item Reason honestly about exactly-once processing on top of at-least-once delivery.
    \item Architect a global telemetry ingestion point for five million mobile-game clients.
\end{itemize}
\end{objectives}

\begin{crosscloud}
Queues, pub/sub, and append-only logs are the three messaging primitives; every cloud ships all three.

\vspace{2mm}
{\footnotesize
\begin{tabular}{@{}p{2.8cm}p{3.0cm}p{3.0cm}p{3.0cm}@{}}
\toprule
\textbf{Capability} & \textbf{AWS} & \textbf{Azure} & \textbf{GCP} \\
\midrule
Pull queue & Amazon SQS & Azure Queue Storage / Service Bus & Pub/Sub (pull) \\
Pub/sub fan-out & Amazon SNS & Event Grid / Service Bus topics & Pub/Sub \\
Durable streaming log & Kinesis Data Streams & Event Hubs & Pub/Sub (ordered) \\
Managed Kafka & Amazon MSK & Event Hubs for Kafka & Managed Service for Apache Kafka \\
Serverless stream processor & Managed Flink / Lambda & Stream Analytics & Dataflow \\
\addlinespace
Scaling unit & Shard (1 MB/s in, 2 MB/s out) & Throughput units / partitions & Throughput quotas \\
Ordering scope & Per partition key & Per partition key & Per ordering key \\
Retention & 24 h--365 days & 1--90 days & Up to 31 days \\
\bottomrule
\end{tabular}}

\vspace{1mm}
\textbf{Fabric} Real-Time Intelligence wraps KQL databases and Eventstreams around the same log concept. \textbf{Snowflake} consumes streams via Snowpipe Streaming. \textbf{MongoDB} change streams are the database-native analog---a log of changes, replayable and resumable.
\end{crosscloud}

\section{Week 13 Roadmap and Mental Model}
Everything so far moved data in files or tables. Week~13 moves it as \emph{events}: small, ordered, continuous, and consumed the moment they exist. The week's core distinction is between a \textbf{queue}---a work-distribution channel where each message is processed once and deleted---and a \textbf{stream}---a durable, replayable log where many consumers read the same events at their own pace \cite{awsKinesisDocs}.

Three services cover the spectrum. \textbf{SQS} is the queue: decoupling, buffering, and per-message deletion. \textbf{SNS} is the fan-out: one publish, many subscribers. \textbf{Kinesis Data Streams (KDS)} is the log: shards, ordering keys, retention, and replay. And \textbf{Amazon MSK} is the ecosystem answer for organizations standardized on Apache Kafka. Choosing among them is Wednesday's work; building correct producers and consumers on the log is Thursday's.

\begin{definitionbox}
\textbf{A streaming data platform} durably records an ordered, append-only sequence of events (the log), lets multiple independent consumers read it at their own pace, and guarantees per-key ordering and at-least-once delivery. Throughput is provisioned in shards (Kinesis) or partitions (Kafka).
\end{definitionbox}

\begin{keyterms}
\textbf{Shard}: Kinesis throughput unit---1 MB/s or 1{,}000 records/s in, 2 MB/s out. \textbf{Partition key}: the hashing input that routes a record to a shard and defines the ordering scope. \textbf{Sequence number}: a record's position in its shard. \textbf{Iterator}: a consumer's cursor into a shard. \textbf{IteratorAgeMilliseconds}: how far behind the newest record a consumer is. \textbf{DLQ}: dead-letter queue for poison messages. \textbf{Fan-out}: one event delivered to many independent consumers (SNS, or enhanced fan-out in KDS).
\end{keyterms}

\section{Monday: Decoupling with SQS and SNS}
\subsection{SQS: The Work Queue}
Amazon SQS is a fully managed queue: producers send messages, consumers poll and delete them, and a visibility timeout hides in-flight messages from other consumers until processing completes or fails \cite{awsSQSDocs}. Standard queues offer at-least-once delivery and best-effort ordering at virtually unlimited throughput; FIFO queues offer exactly-once \emph{delivery} within a deduplication window and strict per-group ordering, capped by throughput quotas per message group. Dead-letter queues capture messages that fail repeatedly, so one poison message cannot block a fleet.
\index{SQS}\index{visibility timeout}\index{dead-letter queue}

\subsection{SNS: The Fan-Out}
Amazon SNS is push-based pub/sub: publishers send to a topic, and every subscription---SQS queues, Lambda functions, HTTPS endpoints, email---receives each message \cite{awsSNSDocs}. The canonical pattern is \emph{fan-out to queues}: one SNS topic fronting several SQS queues, each queue owned by an independent consumer team. SNS handles distribution; SQS gives each consumer its own buffering, retry, and DLQ semantics. Message filtering lets subscribers receive only the subset they declare.
\index{SNS}\index{fan-out pattern}

\begin{tipbox}
Never subscribe a Lambda directly to a high-volume topic when you may need replay or backpressure control. SNS-to-SQS-to-consumer preserves a durable buffer you can drain, inspect, and replay; a direct subscription drops events your consumer could not keep up with.
\end{tipbox}

\section{Tuesday: Kinesis Data Streams---Shards, Keys, and Scaling}
\subsection{The Shard Model}
A Kinesis stream is a sequence of shards; each shard is an ordered, append-only log with fixed capacity: 1 MB/s or 1{,}000 records per second of ingress, 2 MB/s of egress per consumer (or a dedicated 2 MB/s per consumer with enhanced fan-out) \cite{awsKinesisDocs}. Records carry a \emph{partition key}; Kinesis hashes it to a shard, so all records with the same key land in the same shard in order. Retention runs from 24 hours to 365 days, during which any consumer can re-read---the property that makes streams replayable and queues not.
\index{shard}\index{partition key}\index{enhanced fan-out}

\subsection{Sizing and Hot Shards}
Sizing is arithmetic. Ingress: expected events/second $\times$ average record size, divided by 1 MB/s and by 1{,}000 records/s; the shard count is the maximum of the two, plus headroom for growth and skew. The subtle failure is the \emph{hot shard}: if one partition key dominates---a celebrity user, a single tenant---its shard saturates while others idle, and you get \texttt{ProvisionedThroughputExceeded} on a stream that is ``mostly empty.'' Defenses: choose high-cardinality keys, add entropy to hot keys (and handle the ordering consequences), and scale by split/merge or on-demand mode.
\index{hot shard}\index{ProvisionedThroughputExceeded}

\begin{pitfall}
\textbf{Partition keys with low cardinality or skew.} Ordering scope and parallelism are the same choice: the partition key. A key with ten values pins your stream to ten hot spots forever. Choose keys that are high-cardinality \emph{and} the unit of ordering you actually need---usually an entity ID like \texttt{sensor-17} or \texttt{user-92341}.
\end{pitfall}

\subsection{Consumers, Checkpointing, and Iterator Age}
Consumers read via shard iterators and checkpoint their position---the Kinesis Client Library (KCL) does this in DynamoDB, tracking which sequence numbers each worker has processed. The health metric of the whole arrangement is \texttt{IteratorAgeMilliseconds}: the age of the last record the consumer read. Rising iterator age means consumers are falling behind; because retention is finite (24 hours by default), falling far enough behind is data loss. Alarm on it.
\index{KCL}\index{checkpointing}\index{IteratorAgeMilliseconds}

\begin{pitfall}
\textbf{Blocking on poison records.} Kinesis delivers at-least-once and in order per shard: a record your consumer cannot process will be redelivered forever, pinning the shard. Design consumers to be idempotent---deduplicate on a record ID (DynamoDB conditional put, Redis)---and route repeatedly failing records to a dead-letter location (SQS DLQ or S3 error prefix) instead of blocking the iterator.
\end{pitfall}

\section{Wednesday: Amazon MSK Versus Kinesis}
\subsection{Same Log, Different Contract}
Amazon MSK is managed Apache Kafka: real brokers, real ZooKeeper/KRaft, the real Kafka API, in your VPC \cite{awsMSKDocs,apacheKafkaProject}. The comparison with Kinesis is not ``which is better'' but ``which operating model you are buying'':

\begin{table}[H]
\centering
\small
\begin{tabular}{@{}p{3.2cm}p{4.5cm}p{4.7cm}@{}}
\toprule
\textbf{Dimension} & \textbf{Kinesis Data Streams} & \textbf{Amazon MSK} \\
\midrule
API & AWS-native, IAM-integrated & Full Kafka protocol and ecosystem \\
Operations & Nearly none; shard or on-demand scaling & Broker fleet you size, patch windows, version upgrades \\
Pricing unit & Shard-hours + payload (or on-demand) & Broker instance-hours + storage \\
Portability & Low (AWS API) & High (standard Kafka clients move anywhere) \\
Retention & 24 h--365 days, per stream & Configurable, disk-bounded per broker \\
Native integrations & Lambda triggers, Firehose, Flink, KCL & Kafka Connect, MirrorMaker, any Kafka client \\
\bottomrule
\end{tabular}
\caption{Kinesis versus MSK: an operating-model decision.}
\label{tab:ch13-kinesis-msk}
\end{table}

Choose Kinesis when the estate is AWS-native, the team is small, and IAM-integrated access matters. Choose MSK when Kafka compatibility is a requirement---existing clients, Connect connectors, ksqlDB, multi-cloud posture---or when retention and partitioning semantics must be Kafka's exactly.

\section{Thursday Hands-on Project: An IoT Temperature Stream}
\index{hands-on project!Kinesis producer and consumer}
Create a Kinesis data stream, write a Python producer streaming mock IoT temperature readings with per-sensor ordering, and a consumer that prints what it reads.

\subsection*{Lab Objectives}
\begin{itemize}
    \item Create a stream and compute the shard count from first principles.
    \item Produce batches with \texttt{put\_records}, handling partial failures with backoff.
    \item Consume with shard iterators and checkpoint mentally (KCL in production).
    \item Watch \texttt{IteratorAgeMilliseconds} move as the consumer lags and catches up.
\end{itemize}

\subsection*{Procedure}
Fifty sensors reporting once per second at roughly 150 bytes each is $\sim$7.5 KB/s and 50 records/s---one shard is comfortably sufficient; we use two to make the sharding visible.

\begin{lstlisting}[language=bash,caption={Creating the stream.},label={lst:ch13-create-stream}]
aws kinesis create-stream --stream-name iot-temperature-stream --shard-count 2
aws kinesis describe-stream-summary --stream-name iot-temperature-stream
\end{lstlisting}

\begin{lstlisting}[language=Python,caption={Boto3 producer with per-sensor partition keys and partial-failure retry.},label={lst:ch13-producer}]
import json, random, time
import boto3

kinesis = boto3.client("kinesis", region_name="us-east-1")
STREAM = "iot-temperature-stream"

def make_reading(sensor_id: int) -> dict:
    return {
        "sensor_id": f"sensor-{sensor_id}",
        "ts": int(time.time()),
        "temp_c": round(random.uniform(15.0, 35.0), 2),
        "battery_pct": random.randint(20, 100),
    }

while True:
    records = []
    for sensor_id in range(1, 51):
        records.append({
            "Data": json.dumps(make_reading(sensor_id)).encode("utf-8"),
            "PartitionKey": f"sensor-{sensor_id}",  # per-sensor ordering
        })
    resp = kinesis.put_records(StreamName=STREAM, Records=records)
    if resp["FailedRecordCount"]:
        # retry only failed records with backoff (ProvisionedThroughputExceeded)
        failed = [r for r, o in zip(records, resp["Records"]) if "ErrorCode" in o]
        time.sleep(0.5)
        kinesis.put_records(StreamName=STREAM, Records=failed)
    time.sleep(1)
\end{lstlisting}

\begin{lstlisting}[language=Python,caption={Minimal shard-iterator consumer (KCL handles this machinery in production).},label={lst:ch13-consumer}]
import json
import boto3

kinesis = boto3.client("kinesis", region_name="us-east-1")
STREAM = "iot-temperature-stream"

shards = kinesis.list_shards(StreamName=STREAM)["Shards"]
iterators = [
    kinesis.get_shard_iterator(
        StreamName=STREAM,
        ShardId=s["ShardId"],
        ShardIteratorType="LATEST",
    )["ShardIterator"]
    for s in shards
]

while True:
    for i, it in enumerate(iterators):
        resp = kinesis.get_records(ShardIterator=it, Limit=100)
        iterators[i] = resp["NextShardIterator"]
        for rec in resp["Records"]:
            reading = json.loads(rec["Data"])
            print(f"shard {i}: {reading['sensor_id']} "
                  f"{reading['temp_c']}C seq={rec['SequenceNumber']}")
\end{lstlisting}

\subsection*{Verification}
\begin{itemize}
    \item Records from one sensor arrive in timestamp order within their shard---the partition-key guarantee made visible.
    \item Halting the consumer and restarting from \texttt{TRIM\_HORIZON} replays retained records---the stream-is-a-log property.
    \item Forcing throughput past the shard limit (raise the sensor count) produces \texttt{ProvisionedThroughputExceeded} and exercises the retry path.
    \item CloudWatch shows \texttt{IteratorAgeMilliseconds} climbing while the consumer sleeps and falling when it resumes.
\end{itemize}

\section{Friday Use Case: Global Telemetry for Five Million Mobile Games}
\index{use case!mobile game telemetry}
A games publisher must ingest telemetry from five million concurrent mobile clients: session events, purchases, and crash reports, feeding live dashboards, anti-cheat analytics, and a data lake. The requirements state ``exactly-once processing and strict message ordering.'' Your job includes negotiating those requirements into something buildable.

\subsection{Unpacking the Requirements}
\textbf{Ordering} is only definable per entity: events for player \texttt{p-123} must be ordered; player A's events against player B's need no order. That is exactly the partition-key guarantee---so ordering costs nothing if scoped per player, and is impossible (or ruinously serializing) if scoped globally.

\textbf{Exactly-once} splits into delivery and effect. Kinesis delivers at-least-once; exactly-once \emph{effects} come from idempotent consumers: deduplicate on a unique event ID (the client generates a UUID per event; the consumer does a DynamoDB conditional put before writing), or make downstream writes naturally idempotent (upsert by event ID). Where the sink supports transactions (Chapter~15's two-phase commit), end-to-end exactly-once is real; everywhere else, ``exactly-once'' is a contract about effects, not wires.

\subsection{Architecture}
\begin{figure}[H]
    \centering
    \resizebox{\textwidth}{!}{%
    \begin{tikzpicture}[
        box/.style={rectangle, rounded corners, draw=primary, thick, fill=primary!6, align=center, minimum width=2.9cm, minimum height=1.0cm, font=\small\bfseries},
        guard/.style={rectangle, rounded corners, draw=red!60!black, thick, fill=red!5, align=center, minimum width=2.6cm, minimum height=0.9cm, font=\footnotesize\bfseries},
        arr/.style={-{Stealth[length=2.5mm]}, draw=primary, thick}
    ]
        \node[box] (clients) at (0,0) {5M clients\\(SDK batching)};
        \node[box] (edge) at (3.3,0) {API edge\\(validate, auth)};
        \node[box] (kds) at (6.6,0) {Kinesis Data Streams\\on-demand mode};
        \node[box] (flink) at (9.9,1.0) {Flink analytics\\(Ch.~15)};
        \node[box] (fh) at (9.9,-1.0) {Firehose $\rightarrow$ S3\\(Ch.~14)};
        \node[guard] (dedup) at (6.6,2.0) {Idempotent consumers\\event-ID dedup + DLQ};
        \node[guard] (lag) at (9.9,-2.3) {Alarms: IteratorAge,\\WriteThroughputExceeded};
        \draw[arr] (clients) -- (edge);
        \draw[arr] (edge) -- (kds);
        \draw[arr] (kds) -- (flink);
        \draw[arr] (kds) -- (fh);
        \draw[arr, dashed] (dedup) -- (kds);
    \end{tikzpicture}%
    }
    \caption{Global telemetry: edge validation, per-player partition keys, on-demand shards, independent consumers, idempotent writes.}
    \label{fig:ch13-telemetry}
\end{figure}

\subsection{Design Decisions}
\textbf{Client discipline.} Clients batch events (10--50 per request) and tag each with a UUID and monotonically increasing per-session sequence. Batching converts 5M $\times$ 1 event/s into a manageable record rate; UUIDs make dedup possible.

\textbf{Capacity.} Peak estimate: 200{,}000 records/second $\times$ 300 bytes $\approx$ 60 MB/s ingress $\Rightarrow$ roughly 200 shards with headroom (records/s bound). On-demand capacity mode absorbs the daily player-count swings; alarms on write-throughput rejections catch the edge cases.

\textbf{Regional strategy.} Streams are regional; clients write to their nearest region's stream, and cross-region replication into the analytics region consolidates. Global ordering is neither needed nor attempted; per-player ordering survives because a player is pinned to a region per session.

\begin{pitfall}
\textbf{Accepting ``exactly-once, globally ordered'' as a literal requirement.} Global order destroys parallelism; exactly-once delivery does not exist on unreliable networks---only exactly-once effects do. The professional move is to negotiate to per-entity ordering plus idempotent consumers, and write the dedup contract into the API.
\end{pitfall}

\section{Interview Q\&A}
\subsection{Concepts and Trade-offs}
\begin{enumerate}
    \item \textbf{Q: Which two throughput limits determine how many shards a stream needs?}\\
    A: Per-shard ingress of 1 MB/s \emph{and} 1{,}000 records/s---the shard count is the max of the byte-based and record-based computations, plus headroom.
    \item \textbf{Q: Why does partition-key choice control per-entity ordering?}\\
    A: Kinesis hashes the key to a shard; a shard is the ordered unit. Equal keys share a shard, so per-key order is guaranteed---and a skewed key concentrates load on one shard.
    \item \textbf{Q: What does a rising \texttt{IteratorAgeMilliseconds} warn about?}\\
    A: Consumers falling behind. Since retention is finite (default 24 h), a consumer that stays behind long enough loses data; alarm before the retention horizon.
    \item \textbf{Q: How do you build exactly-once effects on at-least-once delivery?}\\
    A: Idempotent consumers: deduplicate on a unique event ID (conditional put in DynamoDB) or write idempotently (upsert by key), and route poison records to a DLQ.
    \item \textbf{Q: Kinesis versus MSK: which is billed per shard-hour and record volume?}\\
    A: Kinesis. MSK bills broker instance-hours and storage---you pay for a fleet; Kinesis bills for log capacity.
    \item \textbf{Q: When does SQS FIFO suffice and Kinesis is overkill?}\\
    A: When each message has exactly one consumer, no replay is needed, and throughput fits FIFO quotas. Streams earn their cost when multiple consumers read the same events or replay matters.
\end{enumerate}

\section{Further Reading}
The Kinesis Data Streams Developer Guide covers shards, partition keys, iterators, KCL checkpointing, and on-demand mode \cite{awsKinesisDocs}. The SQS and SNS guides cover visibility timeouts, DLQs, FIFO, and fan-out \cite{awsSQSDocs,awsSNSDocs}. The MSK Developer Guide and the Kafka documentation cover the broker-side model \cite{awsMSKDocs,apacheKafkaProject}. Streaming Systems provides the delivery-semantics and ordering theory behind Friday's negotiation \cite{akidauStreamingSystems}.

\section*{Key Takeaways}
\begin{itemize}
    \item Queues distribute work (consumed once); streams record a replayable log (read by many). Pick per consumer topology, not habit.
    \item Shard sizing is byte-and-record arithmetic; hot shards are a key-design failure, not a capacity failure.
    \item Partition keys simultaneously define ordering scope and parallelism---choose them as high-cardinality entity IDs.
    \item \texttt{IteratorAgeMilliseconds} against the retention window is the stream's existential alarm.
    \item At-least-once delivery plus idempotent consumers plus DLQs is how exactly-once \emph{effects} are actually built.
    \item Kinesis buys AWS-native simplicity; MSK buys Kafka compatibility and portability.
\end{itemize}

\section{Exercises}
\begin{enumerate}
    \item \textbf{(Conceptual)} Explain the difference between message deletion in SQS and retention in Kinesis, and what each implies for adding a new consumer next month.
    \item \textbf{(Conceptual)} Why does enhanced fan-out exist, given shards already have 2 MB/s of egress?
    \item \textbf{(Hands-on)} Run the Thursday lab; then set the consumer's iterator to \texttt{TRIM\_HORIZON} after ten minutes and document what replays.
    \item \textbf{(Hands-on)} Force a hot shard: send 80\% of records with one partition key and observe which metrics move.
    \item \textbf{(Hands-on)} Build the SNS-to-SQS fan-out: one topic, two queues, message filtering per queue, and a DLQ with a redrive policy.
    \item \textbf{(Design)} Size the shard count for 50{,}000 events/second at 2 KB average, with 3$\times$ burst headroom. Show the arithmetic both ways.
    \item \textbf{(Design)} Design the dedup contract for the Friday telemetry case: where the UUID is generated, where dedup state lives, and what the eviction window must cover.
    \item \textbf{(Architecture)} The publisher must now serve European players from EU-resident infrastructure. Sketch the regional design and the consolidation path.
\end{enumerate}

\section{Capstone Project: A Resilient Telemetry Ingestion Tier}\label{sec:ch13-capstone}
\index{capstone project}\index{Kinesis!capstone}

\begin{capstonebox}
\textbf{Mission.} Build a production-shaped telemetry ingestion tier: a sized and alarmed Kinesis stream, a retrying producer, an idempotent consumer with checkpointing and a DLQ, and the evidence that ordering, replay, and failure semantics all behave as designed.

\textbf{Estimated time.} 4--6 hours in a sandbox AWS account.

\textbf{What you\textquotesingle{}ll practice.}
\begin{itemize}
    \item Shard sizing arithmetic and on-demand versus provisioned mode selection.
    \item Producer partial-failure handling and partition-key design.
    \item Idempotent consumption with dedup state and poison-record routing.
    \item Stream observability: iterator age, throughput rejections, DLQ depth.
\end{itemize}
\end{capstonebox}

\subsection*{Scenario}
The publisher from the Friday use case is piloting with 50{,}000 concurrent clients. Build the ingestion tier they will scale from, and prove its failure semantics before the pilot---not during the launch spike.

\subsection*{Requirements}
\textbf{Functional requirements.}
\begin{itemize}
    \item A stream sized from a documented workload estimate, with the arithmetic committed.
    \item A producer handling partial \texttt{put\_records} failures with bounded backoff.
    \item A consumer that deduplicates on event ID, checkpoints progress, and routes three-time failures to a DLQ.
    \item Alarms on iterator age, write-throughput rejections, and DLQ depth, wired to SNS.
\end{itemize}

\textbf{Non-functional requirements.}
\begin{itemize}
    \item A killed consumer resumes from its checkpoint with no duplicates downstream.
    \item A hot-key fault injection is detectable from metrics alone.
    \item The DLQ is inspectable and a redrive procedure is documented.
\end{itemize}

\subsection*{Implementation Steps}
\begin{enumerate}
    \item Write the sizing memo and create the stream accordingly.
    \item Implement the producer (extend \cref{lst:ch13-producer} with event UUIDs).
    \item Implement the idempotent consumer with DynamoDB conditional-put dedup and checkpointing.
    \item Inject faults: consumer kill, poison record, hot key; capture metric and DLQ evidence for each.
    \item Configure the three alarms and demonstrate one firing.
    \item Write the redrive runbook and the pilot go/no-go checklist.
\end{enumerate}

\subsection*{Deliverables}
\begin{itemize}
    \item The sizing memo, stream configuration, and alarm definitions.
    \item Producer and consumer code with the dedup and DLQ paths.
    \item Fault-injection evidence: kill/resume without duplicates, poison record in the DLQ, hot-key metrics.
    \item The redrive runbook and go/no-go checklist.
\end{itemize}

\subsection*{Acceptance Criteria}
\begin{itemize}
    \item Per-player ordering holds under the full test load, verified programmatically.
    \item Replay from \texttt{TRIM\_HORIZON} produces zero duplicate downstream effects.
    \item Poison records never block a shard longer than the retry policy allows.
    \item All three alarms have fired at least once during testing and are documented.
\end{itemize}

\subsection*{Stretch Goals}
\begin{itemize}
    \item Add a Lambda consumer triggered by the stream with a bisect-batch-on-error policy and compare its failure behavior.
    \item Attach a Firehose delivery stream as a second consumer (preview of Chapter~14).
    \item Replatform the consumer onto KCL with multi-instance rebalancing and document the lease table.
    \item Repeat the pilot sizing for MSK with the same workload and compare monthly cost.
\end{itemize}
